\subsection*{Operand and Status Notation}
Rn names an (:term:base.terms.rn_register:) selected from the (:term:base.terms.general_purpose_register|plural:) R0 through R15; SP and PC are special registers
outside that encoding namespace. An
\texttt{<ea>} operand names a compact effective-address (:term:base.terms.field:) and any appended EA (:term:base.terms.payload:).

FLAGS and STATUS are (:term:base.terms.architectural_state:) registers. Instruction descriptions state whether an
encoding updates, preserves, reads, or ignores each status register. A FLAGS update always names and replaces the
complete ZNCV image.

In a transposed flag-effects table, \texttt{-} marks an unchanged flag and \texttt{0} or \texttt{1} marks a
constant result. \texttt{*} marks the single nonconstant effect defined by that table\textquotesingle{}s local legend. When a table
contains multiple nonconstant effects, letters \texttt{a} through \texttt{z} distinguish their local legend entries.

In instruction pseudocode, Operand Read applies the addressing mode and operation size. Operand Write stores the
selected-size result; an Rn write defines all 64 bits and extends a sub-Q result according to the common integer
result-extension rule, while a memory write changes only the selected bytes. EA Address
computes an address without reading memory. A control-target memory operand reads its declared-width target, while a
control-target register or immediate operand supplies the target value directly. A selector may name an (:term:base.terms.rn_register:), a
segment register, or an encoded immediate. FLAGS ZNCV refers to the integer (:ref:base.registers.STATE.FLAGS.Z:), (:ref:base.registers.STATE.FLAGS.N:), (:ref:base.registers.STATE.FLAGS.C:), and (:ref:base.registers.STATE.FLAGS.V:) bits.
